\documentclass[12pt, twoside]{article}
\input{../preamble}
\usepackage{pgfplots}
\pgfplotsset{compat=1.18}

\fancyhead[LO]{La Scuola d'Italia / Dr.\ Huson / IB Math \\* 27 March 2026}
\fancyhead[RO]{Markscheme}
\fancyfoot[C]{\thepage}

\begin{document}
\subsubsection*{6.2 Classwork: Quadratics --- Markscheme}

\begin{enumerate}[itemsep=0.2cm]

%--- Problem 1: Tangent line to parabola / discriminant [10 marks] ---
\item The diagram below shows the parabola $y = x^2 - 2x$ and the line $y = 4x + k$,
where $k$ is a constant. The line is tangent to the parabola at the point $P$.

\begin{enumerate}[itemsep=0.3cm]
  \item Write down an equation whose solutions give the $x$-coordinates of
        the points where the line meets the parabola. \hfill [2]

  \begin{minipage}[t]{0.62\textwidth}
  $x^2 - 2x = 4x + k$\\[0.15cm]
  $x^2 - 6x - k = 0$
  \end{minipage}%
  \begin{minipage}[t]{0.35\textwidth}\raggedleft
  \textbf{(M1)}\\[0.15cm]
  \textbf{(A1)}
  \end{minipage}

  \medskip
  \item Simplify your equation to the form $x^2 + bx + c = 0$, stating $b$
        and $c$ in terms of $k$. \hfill [2]

  \begin{minipage}[t]{0.62\textwidth}
  $x^2 - 6x + (-k) = 0$\\[0.15cm]
  $b = -6$, \quad $c = -k$
  \end{minipage}%
  \begin{minipage}[t]{0.35\textwidth}\raggedleft
  \textbf{(A1)}\\[0.15cm]
  \textbf{(A1)}
  \end{minipage}

  \medskip
  \item For the line to be tangent, state the condition on the discriminant. \hfill [1]

  \begin{minipage}[t]{0.62\textwidth}
  $\Delta = 0$ \quad (exactly one solution)
  \end{minipage}%
  \begin{minipage}[t]{0.35\textwidth}\raggedleft
  \textbf{(A1)}
  \end{minipage}

  \medskip
  \item Find the value of $k$. \hfill [3]

  \begin{minipage}[t]{0.62\textwidth}
  $\Delta = (-6)^2 - 4(1)(-k) = 36 + 4k$\\[0.15cm]
  $36 + 4k = 0$\\[0.15cm]
  $k = -9$
  \end{minipage}%
  \begin{minipage}[t]{0.35\textwidth}\raggedleft
  \textbf{(M1)}\\[0.15cm]
  \textbf{(A1)}\\[0.15cm]
  \textbf{(A1)}
  \end{minipage}

  \medskip
  \item Find the coordinates of the point $P$. \hfill [2]

  \begin{minipage}[t]{0.62\textwidth}
  $x^2 - 6x + 9 = 0 \Rightarrow (x - 3)^2 = 0$\\[0.15cm]
  $x = 3$, \quad $y = 9 - 6 = 3$\\[0.15cm]
  $P = (3,\; 3)$
  \end{minipage}%
  \begin{minipage}[t]{0.35\textwidth}\raggedleft
  \textbf{(M1)}\\[0.15cm]
  ~\\[0.15cm]
  \textbf{(A1)}
  \end{minipage}
\end{enumerate}

\newpage

%--- Problem 2: Vertex form and graph properties [8 marks] ---
\item Let $f(x) = -x^2 + 4x + 5$.

\begin{enumerate}[itemsep=0.3cm]
  \item Write $f(x)$ in the form $-(x - h)^2 + k$. \hfill [3]

  \begin{minipage}[t]{0.62\textwidth}
  $f(x) = -(x^2 - 4x) + 5$\\[0.15cm]
  $= -(x^2 - 4x + 4) + 4 + 5$\\[0.15cm]
  $= -(x - 2)^2 + 9$
  \end{minipage}%
  \begin{minipage}[t]{0.35\textwidth}\raggedleft
  \textbf{(M1)}\\[0.15cm]
  \textbf{(A1)}\\[0.15cm]
  \textbf{(A1)}
  \end{minipage}

  \medskip
  \item Write down the coordinates of the vertex of the parabola. \hfill [1]

  \begin{minipage}[t]{0.62\textwidth}
  Vertex $= (2,\; 9)$
  \end{minipage}%
  \begin{minipage}[t]{0.35\textwidth}\raggedleft
  \textbf{(A1)}
  \end{minipage}

  \medskip
  \item Write down the equation of the axis of symmetry. \hfill [1]

  \begin{minipage}[t]{0.62\textwidth}
  $x = 2$
  \end{minipage}%
  \begin{minipage}[t]{0.35\textwidth}\raggedleft
  \textbf{(A1)}
  \end{minipage}

  \medskip
  \item Find the $x$-intercepts of the graph of $f$. \hfill [2]

  \begin{minipage}[t]{0.62\textwidth}
  $-(x - 2)^2 + 9 = 0$\\[0.15cm]
  $(x - 2)^2 = 9 \Rightarrow x - 2 = \pm 3$\\[0.15cm]
  $x = -1$ and $x = 5$
  \end{minipage}%
  \begin{minipage}[t]{0.35\textwidth}\raggedleft
  \textbf{(M1)}\\[0.15cm]
  ~\\[0.15cm]
  \textbf{(A1)}
  \end{minipage}

  \medskip
  \item Write down the $y$-intercept. \hfill [1]

  \begin{minipage}[t]{0.62\textwidth}
  $f(0) = 0 + 0 + 5 = 5$\\[0.15cm]
  $y$-intercept $= (0,\; 5)$
  \end{minipage}%
  \begin{minipage}[t]{0.35\textwidth}\raggedleft
  \textbf{(A1)}
  \end{minipage}
\end{enumerate}

\newpage

%--- Problem 3: Discriminant and parameter conditions [8 marks] ---
\item The equation $2x^2 - 4x + p = 0$ has two distinct real roots, where $p \in \mathbb{R}$.

\begin{enumerate}[itemsep=0.3cm]
  \item Find the discriminant of the equation in terms of $p$. \hfill [2]

  \begin{minipage}[t]{0.62\textwidth}
  $\Delta = (-4)^2 - 4(2)(p)$\\[0.15cm]
  $= 16 - 8p$
  \end{minipage}%
  \begin{minipage}[t]{0.35\textwidth}\raggedleft
  \textbf{(M1)}\\[0.15cm]
  \textbf{(A1)}
  \end{minipage}

  \medskip
  \item Find the values of $p$ for which the equation has two distinct real roots. \hfill [2]

  \begin{minipage}[t]{0.62\textwidth}
  $\Delta > 0$\\[0.15cm]
  $16 - 8p > 0 \Rightarrow p < 2$
  \end{minipage}%
  \begin{minipage}[t]{0.35\textwidth}\raggedleft
  \textbf{(M1)}\\[0.15cm]
  \textbf{(A1)}
  \end{minipage}

  \medskip
  \item Find the values of $p$ for which the equation has no real roots. \hfill [1]

  \begin{minipage}[t]{0.62\textwidth}
  $\Delta < 0 \Rightarrow p > 2$
  \end{minipage}%
  \begin{minipage}[t]{0.35\textwidth}\raggedleft
  \textbf{(A1)}
  \end{minipage}

  \medskip
  \item Find the value of $p$ for which there is exactly one real root,
        and find that root. \hfill [3]

  \begin{minipage}[t]{0.62\textwidth}
  $\Delta = 0 \Rightarrow 16 - 8p = 0 \Rightarrow p = 2$\\[0.15cm]
  $2x^2 - 4x + 2 = 0$\\[0.15cm]
  $x^2 - 2x + 1 = 0 \Rightarrow (x - 1)^2 = 0$\\[0.15cm]
  $x = 1$
  \end{minipage}%
  \begin{minipage}[t]{0.35\textwidth}\raggedleft
  \textbf{(A1)}\\[0.15cm]
  \textbf{(M1)}\\[0.15cm]
  ~\\[0.15cm]
  \textbf{(A1)}
  \end{minipage}
\end{enumerate}

\newpage

%--- Problem 4: Line and parabola intersection [9 marks] ---
\item The parabola $y = x^2 - 2x - 3$ and the line $y = x + 1$ are drawn
on the same axes.

\begin{enumerate}[itemsep=0.3cm]
  \item Show that the $x$-coordinates of the points of intersection satisfy\\
        $x^2 - 3x - 4 = 0$. \hfill [2]

  \begin{minipage}[t]{0.62\textwidth}
  $x^2 - 2x - 3 = x + 1$\\[0.15cm]
  $x^2 - 3x - 4 = 0$
  \end{minipage}%
  \begin{minipage}[t]{0.35\textwidth}\raggedleft
  \textbf{(M1)}\\[0.15cm]
  \textbf{(A1)} \textbf{(AG)}
  \end{minipage}

  \medskip
  \item Factorise $x^2 - 3x - 4$. \hfill [1]

  \begin{minipage}[t]{0.62\textwidth}
  $(x - 4)(x + 1)$
  \end{minipage}%
  \begin{minipage}[t]{0.35\textwidth}\raggedleft
  \textbf{(A1)}
  \end{minipage}

  \medskip
  \item Hence find the coordinates of the two points of intersection. \hfill [3]

  \begin{minipage}[t]{0.62\textwidth}
  $x = 4$: \quad $y = 4 + 1 = 5$ \quad $\Rightarrow (4,\; 5)$\\[0.15cm]
  $x = -1$: \quad $y = -1 + 1 = 0$ \quad $\Rightarrow (-1,\; 0)$
  \end{minipage}%
  \begin{minipage}[t]{0.35\textwidth}\raggedleft
  \textbf{(A1)}\\[0.15cm]
  \textbf{(A1)(A1)}
  \end{minipage}

  \medskip
  \item Find the equation of the axis of symmetry of the parabola. \hfill [1]

  \begin{minipage}[t]{0.62\textwidth}
  $x = \dfrac{-(-2)}{2(1)} = 1$
  \end{minipage}%
  \begin{minipage}[t]{0.35\textwidth}\raggedleft
  \textbf{(A1)}
  \end{minipage}

  \medskip
  \item Find the coordinates of the vertex of the parabola. \hfill [2]

  \begin{minipage}[t]{0.62\textwidth}
  $f(1) = 1 - 2 - 3 = -4$\\[0.15cm]
  Vertex $= (1,\; -4)$
  \end{minipage}%
  \begin{minipage}[t]{0.35\textwidth}\raggedleft
  \textbf{(M1)}\\[0.15cm]
  \textbf{(A1)}
  \end{minipage}
\end{enumerate}

\end{enumerate}
\end{document}
